\section{Experimental design}\label{sec:setup}

The two domains were chosen to differ in the \emph{type} of error the agent makes, not only in difficulty. This
contrast is what makes RQ3 testable. Path planning has an exact, microsecond validator, so it isolates the
quality of the screens. The CSTR has a simulation-based validator, so it measures safety and compute in an
engineering setting.

\subsection{Domain 1: graph path planning with an exact validator}
Each instance is a directed graph with 5, 10, 15 or 20 nodes. It is given in the prompt as one adjacency line per
node (``\texttt{u: v1, v2}''), together with a start node, a goal node and one worked example on a different
graph. The agent returns the path as a JSON list. The validator checks that the path starts and ends at the
right nodes and that every step $u\to v$ is an edge, that is, that $v$ appears in the adjacency line of $u$.
Graphs are disjoint across all data sets. For each model, probes were developed on 3000 graphs (1500 training,
500 calibration, 500 threshold selection, 500 test). They were evaluated once on a fresh certification set of
3000 new graphs that had not been used in any earlier analysis. A fresh set was used because the
original test and certification graphs had been examined in earlier exploratory analyses.

\subsection{Domain 2: CSTR fault recovery with a digital-twin validator}
The plant is the CSTR case study of \citet{car_todo} (Fig.~\ref{fig:cstr}). An exothermic reaction
$A\rightarrow B$ takes place in a jacketed vessel. Three PID loops regulate the feed flow (FC), the level, by the
outlet pump (LC), and the temperature, by the coolant valve (TC). A monitoring layer triggers the agent when the
process leaves its safe operating region. Each episode draws a nominal plant: heat-transfer coefficient, feed and
coolant temperatures, inlet concentration, coolant capacity, and nominal level and feed setpoints. It also draws
one of four fault families (jacket fouling, outlet-pump degradation, cooling valve stuck partly closed, outlet
blockage), a continuous fault severity and an onset time between 1500 and 5000\,s. The plant is simulated until
the first monitor trigger, and the snapshot at that moment is the agent's task. Episodes that never trigger, or
that trigger before the fault onset, are redrawn within the same fault family.

\begin{figure}[!htbp]
\centering
% CSTR case study: process, regulatory loops, fault locations and the supervisory agent loop.
\resizebox{\textwidth}{!}{%
\begin{tikzpicture}[font=\footnotesize, >={Latex[length=1.8mm]},
  pipe/.style={line width=1.1pt, black!75},
  sig/.style={densely dashed, black!60},
  ctrl/.style={circle, draw=black!70, fill=white, minimum size=7.5mm, inner sep=0pt, font=\scriptsize},
  fault/.style={circle, draw=red!75!black, fill=red!12, text=red!75!black, minimum size=4.6mm, inner sep=0pt, font=\scriptsize\bfseries},
  sup/.style={draw=black!70, rounded corners=2pt, minimum width=31mm, minimum height=7.5mm, align=center, fill=black!3},
]
\newcommand{\valve}[1]{\fill[white] ($(#1)+(-0.22,-0.15)$) -- ($(#1)+(0.22,0.15)$) -- ($(#1)+(0.22,-0.15)$) -- ($(#1)+(-0.22,0.15)$) -- cycle;
  \draw[black!80, line width=0.7pt] ($(#1)+(-0.22,-0.15)$) -- ($(#1)+(0.22,0.15)$) -- ($(#1)+(0.22,-0.15)$) -- ($(#1)+(-0.22,0.15)$) -- cycle;}

% jacket and vessel
\fill[cTwin!30] (-0.35,-0.35) rectangle (2.55,2.4);
\draw[black!70] (-0.35,-0.35) rectangle (2.55,2.4);
\fill[white] (0,0) rectangle (2.2,3.2);
\fill[cTwin!12] (0,0) rectangle (2.2,2.05);
\draw[black!80, line width=1pt] (0,3.2) -- (0,0) -- (2.2,0) -- (2.2,3.2);
\draw[black!50, decorate, decoration={snake, amplitude=0.6pt, segment length=4mm}] (0,2.05) -- (2.2,2.05);
\node[font=\scriptsize, align=center] at (1.1,1.0) {$A\rightarrow B$\\[-1pt]exothermic};
\node[font=\scriptsize, text=black!70, rotate=90] at (2.38,1.0) {jacket};

% feed line with flow controller
\draw[pipe, ->] (-3.2,3.6) -- (0.6,3.6) -- (0.6,2.3);
\valve{-1.2,3.6}
\node[anchor=south west, font=\scriptsize, align=left] at (-3.2,3.65) {feed $F_{\mathrm{in}}$, $C_{A,\mathrm{in}}$, $T_{\mathrm{in}}$};
\node[ctrl] (fc) at (-1.2,4.55) {FC};
\draw[sig] (fc) -- (-1.2,3.75);

% outlet with pump and level controller
\draw[pipe, ->] (1.1,-0.35) -- (1.1,-1.3) -- (4.4,-1.3);
\node[circle, draw=black!80, fill=white, minimum size=6.5mm, line width=0.7pt] (pump) at (2.9,-1.3) {};
\draw[black!80] (pump.north) -- (pump.east) -- (pump.south);
\node[anchor=north, font=\scriptsize] at (3.95,-1.35) {outlet};
\node[ctrl] (lc) at (3.65,1.75) {LC};
\draw[sig] (2.2,1.75) -- node[above, font=\tiny] {LT} (lc);
\draw[sig] (lc) -- (3.65,-0.25) -| (pump.north);

% cooling jacket with temperature controller
\draw[pipe, ->] (-3.2,0.15) -- (-0.35,0.15);
\draw[pipe, ->] (-0.35,2.0) -- (-3.2,2.0);
\valve{-1.9,0.15}
\node[anchor=north west, font=\scriptsize] at (-3.2,0.1) {coolant in};
\node[anchor=south west, font=\scriptsize] at (-3.2,2.05) {coolant out};
\node[ctrl] (tc) at (-1.9,1.1) {TC};
\draw[sig] (tc) -- (-1.9,0.3);
\draw[sig] (tc) -- node[above, font=\tiny, pos=0.75] {TT} (0.45,1.1);

% faults
\node[fault] at (-0.17,0.55) {1};
\node[fault] at (2.9,-0.8) {2};
\node[fault] at (-1.4,-0.2) {3};
\node[fault] at (1.1,-0.85) {4};
\node[anchor=north west, font=\scriptsize, align=left, text=red!60!black] at (-3.25,-1.75)
  {faults: \textbf{1} jacket fouling (lower $UA$)\quad \textbf{2} outlet pump degradation\\
   \phantom{faults: }\textbf{3} cooling valve stuck partly closed\quad \textbf{4} outlet blockage};

% ---------------------------------------------------------------- supervisory layer
\node[sup] (mon) at (7.6,4.4) {monitor: process leaves\\[-1pt]safe operating region};
\node[sup] (agent) at (7.6,3.15) {LLM agent proposes\\[-1pt]$(T_{\mathrm{sp}},L_{\mathrm{sp}},F_{\mathrm{in,sp}})$};
\node[sup, draw=cInt, fill=cInt!8] (screen) at (7.6,1.9) {risk screen on\\[-1pt]internals (+ observables)};
\node[sup, draw=cTwin!90!black, fill=cTwin!15] (twin) at (7.6,0.65) {digital twin: rollout to\\[-1pt]600\,s horizon};
\node[sup] (sp) at (7.6,-0.6) {apply setpoints};
\draw[->] (mon) -- (agent);
\draw[->] (agent) -- (screen);
\draw[->] (screen) -- node[right, font=\scriptsize] {\VAL} (twin);
\draw[->] (twin) -- node[right, font=\scriptsize] {pass} (sp);
\draw[->] (screen.east) -- ++(0.45,0) |- node[pos=0.25, right, font=\scriptsize] {\ACC} (sp.east);
\draw[->] (screen.west) -- ++(-0.35,0) |- node[pos=0.25, left, font=\scriptsize] {\REJ} (agent.west);
\draw[->] (twin.west) -- ++(-0.55,0) |- node[pos=0.2, left, font=\scriptsize] {fail} (agent.west);
% measurements up, setpoints down
\draw[->, sig, black!70] (2.4,3.2) -- node[above, font=\scriptsize, sloped] {plant snapshot} (mon.west);
\draw[->, line width=0.8pt, cRzero!90!black] (sp.west) -- (4.5,-0.6);
\node[font=\scriptsize, text=cRzero!80!black, align=center, anchor=south] at (5.2,-0.55) {setpoints to\\[-1pt]TC, LC, FC};
\end{tikzpicture}}

\caption{CSTR case study. Left: jacketed reactor with flow (FC), level (LC) and temperature (TC) loops and the
four fault locations. Right: supervisory loop. When the monitor detects that the process has left its safe
operating region, the LLM agent proposes new setpoints for the three loops. The risk screen accepts the proposal
without simulation, rejects it without simulation, or passes it to the digital twin, which simulates the plant
forward from the snapshot.}
\label{fig:cstr}
\end{figure}

The prompt contains the agent's role and the plant's knowledge graph in Turtle syntax: process operators,
controllers, and the mass and energy balance equations. It also contains the snapshot (measurements, actuator
commands, alarms and current setpoints), the goal and the admissible setpoint ranges. The prompt is deliberately
\emph{neutral}: it states the plant physics and the acceptance criteria but gives no advice on which setpoint to
change, in which direction or by how much. Failures therefore reflect the model's own reasoning rather than the
wording of the prompt. The agent returns three setpoints ($T_{\mathrm{sp}}$, $L_{\mathrm{sp}}$,
$F_{\mathrm{in,sp}}$) as JSON after a short justification.

The validator simulates the plant forward from the snapshot under the proposed setpoints, with the plant's own
future disturbances, to the end of the horizon. A proposal passes if the plant reaches the safe state within
600\,s, holds it for 60 consecutive seconds, and remains safe for at least 80\% of the remaining run. The safe
state means a temperature of at most 313\,K, a level between 2 and 11, no overflow, and no actuator saturated
while the temperature is away from 310\,K. One validation took 8--12\,s of wall time on our hardware, with four
parallel simulation workers.

Fault severities were calibrated without the LLM so that the reference model's (Qwen2.5-3B) first proposals pass
in roughly 30\% of episodes. This calibration is a property of the data set, not a target for the agents. It
ensures that enough passing proposals exist to train and certify an \ACC{} decision. First-proposal failure rates
on the evaluation sets range from 66\% to 84\% across models (Table~\ref{tab:datasets}). The 5000 episodes are
partitioned by episode into training (3000), calibration (500), threshold selection (500), test (500) and
certification (500), with equal numbers per fault family in each partition (Fig.~\ref{fig:protocol}). A further
400 fresh episodes, 100 per fault family and disjoint from all others, are reserved for the closed loop.

\begin{figure}[!htbp]
\centering
% Data partitions and their single use; widths are proportional to the number of graphs/episodes.
\resizebox{\textwidth}{!}{%
\begin{tikzpicture}[font=\footnotesize,
  part/.style={draw=white, line width=0.8pt, minimum height=9mm, align=center, font=\scriptsize, inner sep=0pt, anchor=west},
]
\def\u{0.0024} % cm per item
\colorlet{pFit}{black!22}\colorlet{pCal}{cTwin!45}\colorlet{pThr}{cRzero!55}\colorlet{pTest}{cInt!45}
\colorlet{pCert}{cInt!85}\colorlet{pCL}{cGrd!55}
% path planning
\node[anchor=east, font=\footnotesize\bfseries, align=right] at (-0.2,1.5) {Path planning\\[-1pt]{\scriptsize\mdseries per model}};
\node[part, fill=pFit, minimum width=1500*\u cm] (a1) at (0,1.5) {train\\1500};
\node[part, fill=pCal, minimum width=500*\u cm] (a2) at (a1.east) {calib.\\500};
\node[part, fill=pThr, minimum width=500*\u cm] (a3) at (a2.east) {thresh.\\500};
\node[part, fill=pTest, minimum width=500*\u cm] (a4) at (a3.east) {test\\500};
\node[part, fill=pCert, text=white, minimum width=3000*\u cm] (a5) at ($(a4.east)+(0.25,0)$) {fresh certification set\\3000 new graphs (cert2)};
% CSTR
\node[anchor=east, font=\footnotesize\bfseries, align=right] at (-0.2,0) {CSTR\\[-1pt]{\scriptsize\mdseries per model}};
\node[part, fill=pFit, minimum width=3000*\u cm] (b1) at (0,0) {train\\3000 episodes};
\node[part, fill=pCal, minimum width=500*\u cm] (b2) at (b1.east) {calib.\\500};
\node[part, fill=pThr, minimum width=500*\u cm] (b3) at (b2.east) {thresh.\\500};
\node[part, fill=pTest, minimum width=500*\u cm] (b4) at (b3.east) {test\\500};
\node[part, fill=pCert, text=white, minimum width=500*\u cm] (b5) at (b4.east) {cert.\\500};
\node[part, fill=pCL, minimum width=400*\u cm] (b6) at ($(b5.east)+(0.25,0)$) {closed\\loop\\400};
% legend
\begin{scope}[yshift=-1.25cm, font=\scriptsize]
\foreach \c/\t/\x in {pFit/{fit $\ell_2$-logistic probe}/0, pCal/{Platt calibration}/4.4, pThr/{choose $t_{\mathrm{acc}},t_{\mathrm{rej}}$}/8.8} {
  \fill[\c] (\x,0) rectangle ++(0.3,0.22); \node[anchor=west] at (\x+0.35,0.11) {\t}; }
\foreach \c/\t/\x in {pTest/{sealed: AUROC, once}/0, pCert/{sealed: certification, once}/4.4, pCL/{fresh: closed-loop policies}/8.8} {
  \fill[\c] (\x,-0.4) rectangle ++(0.3,0.22); \node[anchor=west] at (\x+0.35,-0.29) {\t}; }
\end{scope}
\draw[decorate, decoration={brace, amplitude=3pt}, black!60] ($(a2.north west)+(0,0.05)$) -- node[above, font=\scriptsize] {development (pooled for thresholds)} ($(a3.north east)+(0,0.05)$);
\end{tikzpicture}}

\caption{Data partitions per model (widths proportional to size). Each partition has a single role.
Thresholds are chosen on the pooled calibration and threshold-selection partitions. The test and certification
partitions were sealed until the screens were frozen and were then evaluated once. The closed-loop episodes are
fresh.}
\label{fig:protocol}
\end{figure}

\begin{table}[!htbp]
\centering
\small
\caption{Evaluation sets used for the confirmatory results (one candidate per graph or episode). Path planning: fresh certification set, never used before evaluation. CSTR: sealed test and certification partitions, each evaluated once.}
\label{tab:datasets}
\begin{tabular}{llrrr}
\toprule
Domain & Model & Set & Eligible & Failure rate \\
\midrule
Path planning & Qwen-3B & cert2 (fresh) & 2992 & 0.51 \\
Path planning & Llama-3B & cert2 (fresh) & 2995 & 0.63 \\
Path planning & Qwen-1.5B & cert2 (fresh) & 2945 & 0.68 \\
Path planning & SmolLM-1.7B & cert2 (fresh) & 2946 & 0.83 \\
Path planning & Qwen-7B & cert2 (fresh) & 2973 & 0.49 \\
CSTR & Qwen-1.5B & test\_iid & 494 & 0.67 \\
CSTR & Qwen-1.5B & cert & 492 & 0.66 \\
CSTR & Qwen-3B$^\dagger$ & test\_iid & 493 & 0.71 \\
CSTR & Qwen-3B$^\dagger$ & cert & 498 & 0.67 \\
CSTR & Qwen-7B & test\_iid & 499 & 0.79 \\
CSTR & Qwen-7B & cert & 498 & 0.80 \\
CSTR & Llama-3B & test\_iid & 473 & 0.84 \\
CSTR & Llama-3B & cert & 474 & 0.83 \\
CSTR & SmolLM-1.7B & test\_iid & 499 & 0.67 \\
CSTR & SmolLM-1.7B & cert & 496 & 0.66 \\
\bottomrule
\end{tabular}
\end{table}


\subsection{Models and inference}
Qwen2.5-1.5B, Qwen2.5-3B and Qwen2.5-7B \citep{qwen2025}, Llama-3.2-3B \citep{llama3} and SmolLM2-1.7B
\citep{allal2025smollm2}, all instruction-tuned, were run locally in bfloat16. The exception is Qwen2.5-7B,
which was run with 4-bit NF4 weights \citep{dettmers2023qlora} to fit an 8\,GB laptop GPU (NVIDIA RTX~4060).
Decoding is greedy, so every proposal is reproducible. The knowledge graph was exported once and frozen, so
prompts are identical across runs. The models span three families and 1.5 to 7 billion parameters. They
were chosen because their weights, attention and hidden states are accessible. API-served models do not expose
these signals, a limitation we return to in Section~\ref{sec:discussion}.

\subsection{Pre-registration, data hygiene and statistics}
For each study we froze a written pre-registration before computing the relevant features or outcomes. It
fixed the hypotheses, the direction of each predicted effect, the primary features, the decision rules and the
analysis \citep{nosek2018preregistration}. Later changes were recorded as dated amendments made before the
affected data were examined: the change of primary endpoint from AUROC to validator calls skipped, the
bound-based accept rule, two added signal combinations that isolate grounding, the closed-loop policies, the
internals-only closed-loop policy, the first-proposal-only rule, and a pilot gate for SmolLM2 in the CSTR. Test
and certification partitions were evaluated once per model with frozen probes and thresholds.

Two analyses are labelled exploratory. The first is the offline CSTR result for Qwen2.5-3B under the bound-based
rule: its sealed partitions had been used once under the earlier point rule, before the amendment. The second is
the CSTR closed-loop probes of Qwen2.5-3B, which come from that exploratory freeze; the closed-loop episodes
themselves are fresh. Where a run failed for technical reasons (GPU memory exhaustion, a crashed simulator
worker), it was resumed from its saved records, and no records were discarded.

The unit of independence is the graph or episode. Differences between screens on the same proposals are reported
as paired differences with 95\% bootstrap confidence intervals over graphs or episodes (2000 resamples), and a
difference is called significant when its interval excludes zero. Closed-loop policies were run on the same 400
episodes and are compared with always-validate by paired bootstrap over episodes. Certification uses the
Bonferroni-corrected exact bound of Proposition~\ref{prop:cert}. All numbers in the tables and figures are
generated by scripts from the saved results.
